\section{Puesta en marcha}

Esta sección describe cómo reproducir la compilación, el flasheo y las pruebas de la PoC.

\subsection{Requisitos}

\begin{itemize}
  \item \dNcs{} con su cadena de herramientas, instalada con \texttt{nrfutil sdk-manager}.
  \item El comando \texttt{nrfutil device}, que el flasheo usa por dentro (\texttt{nrfutil install device}).
  \item Linux, o Windows con WSL2. En WSL2 hay que compartir la placa con \texttt{usbipd-win}: se
        ejecuta una vez como administrador \texttt{usbipd bind} sobre el identificador de la placa,
        y después \texttt{usbipd attach --wsl} cada vez que se conecta.
  \item Una regla de \texttt{udev} que dé permiso sobre los dispositivos de Nordic
        (fabricante \texttt{1915}); sin ella, el flasheo falla con el error ``The device has no
        langid''.
  \item Una ruta del proyecto \emph{sin espacios}: Zephyr no soporta espacios en la ruta de la
        aplicación. El script de compilación lo resuelve con un enlace simbólico.
\end{itemize}

\subsection{Compilación y flasheo}

Desde la carpeta \texttt{firmware/} se usa un solo script, que compila por el enlace simbólico y
deja las compilaciones de medición en directorios aparte para que nunca pisen la normal:

\begin{table}[H]
\centering
\small
\begin{tabularx}{\textwidth}{|l X|}
\hline
\tablehead{\textbf{Comando} & \textbf{Qué hace}} \\ \hline
\texttt{./build.sh} y \texttt{./build.sh flash} & Compila y flashea el firmware normal por USB (MCUboot). \\ \hline
\texttt{./build.sh bench} y \texttt{flash-bench} & Compila y flashea el firmware con el benchmark de \texttt{geofence\_check()}. \\ \hline
\texttt{./build.sh cold} y \texttt{flash-cold} & Compila y flashea el firmware que borra los datos de satélites en cada arranque, para medir el arranque en frío. \\ \hline
\texttt{./build.sh clean} & Borra los directorios de compilación. \\ \hline
\end{tabularx}
\caption{Comandos de \texttt{build.sh}.}
\end{table}

Las dos variantes de medición se activan con opciones de compilación que vienen apagadas
(\texttt{CONFIG\_DRVF\_GEOFENCE\_BENCHMARK} y \texttt{CONFIG\_DRVF\_GNSS\_COLD\_START}), de modo
que el firmware normal no cambia. El flasheo apenas imprime nada: la forma de comprobar que
funcionó es leer la consola.

\subsection{Consola}

La consola es un puerto serie virtual por USB, a 115\,200 baudios (en Linux, \texttt{/dev/ttyACM0}).
Cada línea empieza con el tiempo desde el último reinicio. Al arrancar aparece la versión del
firmware del módem, y después, cada segundo, los satélites listados por sistema y, con fix, la
posición, la precisión reportada, la distancia al borde, el color y el tiempo de la comprobación.

\subsection{Pruebas en el PC}

Los módulos de C puro tienen sus pruebas en el PC, sin placa. En \texttt{firmware/src/geofence}:
\texttt{make test}, \texttt{make stress} (unos 3.8 millones de casos aleatorios con semilla fija),
\texttt{make sanitize} (las mismas pruebas con detectores de errores de memoria y de comportamiento
indefinido) y \texttt{make cross-check} (compila el motor para Cortex-M33 y ESP32-S3 con las
advertencias como errores). En \texttt{firmware/src/alert}, \texttt{fence}, \texttt{app}, \texttt{demo}
y \texttt{bench}: \texttt{make test}.
